\documentclass[12pt]{article}
\usepackage{resource/shortex}
\usepackage{mathtools,amssymb}
\usepackage{enumitem}
\usepackage{datetime}
\usepackage{fancyhdr}
\usepackage{ifthen}
\usepackage{pifont}
\usepackage{mathrsfs}
\usepackage[normalem]{ulem}
%\RequirePackage[usenames,dvipsnames]{color}


%% showing solutions 
\newboolean{showsols}
\setboolean{showsols}{false}
\newcommand{\showsolutions}{\setboolean{showsols}{true}}
\newlength{\solspace}
\setlength{\solspace}{16em}
\newcommand{\nosolspace}{\setlength{\solspace}{0em}}

\newcommand{\solution}[2]{\ifthenelse{\boolean{showsols}}{{\par\textcolor{red}{Rubric criterion: #1}\newline\textcolor{blue}{#2}}}{\vspace{\solspace}}}
\newcommand{\npforprint}{\ifthenelse{\boolean{showsols}}{}{\newpage}}
\newcommand{\vsforprint}[1]{\ifthenelse{\boolean{showsols}}{}{\vspace{#1}}}

%% grading
\newcommand{\grades}[1]{~\newline\noindent\framebox[\textwidth][c]{E\quad \quad \quad \quad S \quad \quad \quad \quad R \quad \quad \quad \quad N}}

\newcommand{\questiongrades}{\noindent\emph{Each item is marked $\checkmark$, $\checkmark\!-$, or \ding{55}:}\begin{center}\emph{$\checkmark$ = capable  $~~$ $\checkmark\!-$ = mostly capable  $~~$ \ding{55} = not yet}\end{center} }

\newcommand{\triAssessments}{$\checkmark\quad\checkmark\!-\quad$\ding{55}\quad}
\newcommand{\biAssessments}{$\checkmark\quad$\ding{55}\quad}

%% line for writing name and BUID
\newcommand{\nameBUID}{\noindent\textbf{Name:}\underline{\phantom{XXXXXXXXXXXXXXXXXXX}}\hfill\textbf{BUID:}\underline{\phantom{XXXXXXXXXXXXX}}\newline}


%% formatting 
%\allowdisplaybreaks[3]
\renewcommand{\headrulewidth}{.4mm} % header line width
\renewcommand{\arraystretch}{1.25}


%% page style 
\pagestyle{fancy}
\fancyhf{}
\rhead{Spring 2026}
\lhead{CAS MA 214: Applied Statistics}
%\rfoot{\thepage}

\begin{document}

\begin{center}
\textbf{\Large Project I Rubric}
\end{center}

\grades
\subsection*{Rubric}
\questiongrades
\noindent
\section*{General Rubric}

\subsubsection*{Data Description}
\normalsize
\begin{tabular}{lp{.75\textwidth}}
\toprule
\textbf{Mark} & \textbf{Description} \\
\hline
\triAssessments &\textbf{(1)} Description of the dataset was clear and well conveyed, including its source and context. \\
\triAssessments &\textbf{(2)} Data summarization and preprocessing steps were appropriately explained and justified. \\
\triAssessments &\textbf{(3)} Research objectives were clearly stated and aligned with the analysis. \\
\bottomrule
\end{tabular}
\vspace{1em}

\subsubsection*{Assumptions and Diagnostics}
\begin{tabular}{lp{.75\textwidth}}
\toprule
\textbf{Mark} & \textbf{Description} \\
\hline
\triAssessments &\textbf{(1)} Choice of predictor and response variables was reasonable and justified in the research context. \\
\triAssessments &\textbf{(2)} Model assumptions (linearity, normality, and independence) were checked and discussed using appropriate plots or tests. \\
\triAssessments &\textbf{(3)} Diagnostic plots and residual analyses were properly interpreted, and any necessary remedial actions were considered. \\
\bottomrule
\end{tabular}
\vspace{1em}

\subsubsection*{Fitting and Interpretation}
\begin{tabular}{lp{.75\textwidth}}
\toprule
\textbf{Mark} & \textbf{Description} \\
\hline
\triAssessments &\textbf{(1)} Model fitting was performed correctly using appropriate R functions and parameters. \\
\triAssessments &\textbf{(2)} Results were accurately reported and interpreted in context. \\
\triAssessments &\textbf{(3)} Conclusions were consistent with the statistical findings and addressed the stated objectives. \\
\bottomrule
\end{tabular}
\vspace{1em}

\subsubsection*{Structure and R Code}
\begin{tabular}{lp{.75\textwidth}}
\toprule
\textbf{Mark} & \textbf{Description} \\
\hline
\triAssessments &\textbf{(1)} Used clear visuals and a well-structured presentation format. \\
\triAssessments &\textbf{(2)} Delivered a clear and logical presentation with smooth transitions and key takeaways. \\
\triAssessments &\textbf{(3)} Used R code that was clean, reproducible, and properly commented. \\
\bottomrule
\end{tabular}
\vspace{1em}

\section*{Individual Rubric}
\begin{tabular}{lp{.75\textwidth}}
\toprule
\textbf{Mark} & \textbf{Description} \\
\hline
\triAssessments &\textbf{(1)} Successfully completed the individual tasks. \\
\triAssessments &\textbf{(2)} Answered all three reflection questions thoughtfully and accurately. \\
\triAssessments &\textbf{(3)} Asked three relevant and valid questions during other group presentations. \\
\bottomrule
\end{tabular}

\newpage
\subsection*{Grading Standards}
\begin{itemize}
\item \textbf{E: Excellent}: Satisfactory with exceptional creativity or insight
\item \textbf{S: Satisfactory}: Any combination of $\checkmark$ and $\checkmark\!-$
\item \textbf{R: Revisions Needed}: Any \ding{55} on any criterion
\item \textbf{N: Not Assessable}: Cannot grade due to too many missing components
\end{itemize}

\end{document}
